\chapter*{Conclusions}\label{chap:conclusion}
\addcontentsline{toc}{chapter}{Conclusions}

The roomlight dependent measurements show our spectrometer is resolving different wavelight photons very well. Although the roomlight source is highly dominating in our readings, after subtracting the dark counts, the LED signal was still visible when the roomlight was turned on. We did still proceed to turn off the roomlight for a cleaner signal. Our radial symmetry measurements did give us a symmetric picture, however the center of the distribution was shifted by some degrees instead of being dead center at 0 degrees. This is probably due to misaligned spectrometer arms.\\

The final measurement, which is the intensity measurement, gave predictable results of an exponentially decreasing intensity as we increase the distance of excitation from the fiber end. Only surprise was from what we understand are a few kinks in the fiber itself which causes a sudden drop in intensity in some LED movement intervals.\\

When it comes to simulation, we had to remove unphysical data points, however what we removed were always small compared to the whole dataset. The maximum reflection angles were visible in the angular distribution of the simulation data and were in agreement with the calculated ones. The intensity of photons seemed to increase linearly with increased angle until the geometrically allowed limit is reached. This is due to the spherical nature of the photon creation process, photons are radiated isotropically.\\

Very clear results for the geometrical analysis can be observed in \hyperref[corerefl]{Figure 2.7} and \hyperref[cladrefl]{Figure 2.8}. The distributions show that the cladding photons have an allowed range in which they can exist with relation to their minimal distance from the center axis and are cut off at an angle of about $21^{\circ}$. On the other hand, the core photons are found where cladding photons are absent. As the reflection angle increases the $r_{\text {min }}$ value also increases for the photon to still propagate in an allowed trajectory of total reflection. This behavior draws a helical shape path inside the fiber in our theory.\\

The biggest inconsistency with this laboratory report can be observed in the determination of the angular dependence of the fit parameter $a$. For any of the data samples, the fit parameter does not show clear dependence on the reflection angle. Moreover, for the theoretical curve there seems to be no dependence at all. As if that was not confusing enough, we observe that when the theory curve is taken as a symmetry axis, the measurement results and simulation results mirror each other. Although, the error allowance added from the evaluation of the fitting parameters allow the values to have crossover at some points. This is the final result of our report.